\documentclass[11pt]{article}
\usepackage[letterpaper,margin=0.82in]{geometry}
\usepackage{amsmath,amssymb,amsthm,mathtools}
\usepackage{graphicx}
\usepackage{microtype}
\usepackage[hidelinks]{hyperref}
\usepackage{enumitem}

\newtheorem{theorem}{Theorem}
\newtheorem{proposition}{Proposition}
\newcommand{\R}{\mathbb{R}}

\title{A smooth counterexample to the negative-\(p\)\\normalized-moment conjecture}
\author{Automated mathematics run \texttt{fa\_banach\_001}\\Lane 3 --- candidate for human review}
\date{17 August 2026}

\begin{document}
\maketitle

\noindent\textbf{Status.}
Candidate full counterexample, likely valid.  The construction disproves the
negative-exponent extension proposed in Remark~2.2 of Jie Xiao's
\emph{Optimal Monotonicity of \(L^p\) Integral of Conformal Invariant Green
Function}, arXiv:0708.2644~\cite{Xiao}.

\section{Source conjecture}

Theorem~2.1(i) of the source fixes \(c>0\), assumes that \(X\geq0\) on
\((0,\infty)\) and
\[
 X'(t)\leq0,
 \qquad
 \frac{d}{dt}\bigl(e^{ct}X(t)\bigr)\leq0,
\]
and defines
\[
 Y_p(0)=-\int_0^\infty t^p\,dX(t).
\]
For \(0\leq p_1<p_2\), it proves
\[
 \frac{c^{p_2}Y_{p_2}(0)}{\Gamma(p_2+1)}
 \leq
 \frac{c^{p_1}Y_{p_1}(0)}{\Gamma(p_1+1)}.                 \tag{1}
\]
Remark~2.2, on source-PDF page 6, conjectures that Theorem~2.1(i)
remains valid for \(-1<p_1<0\), and explicitly discusses
\(-1<p_1<p_2<0\).  The original passage is reproduced below.

\begin{center}
 \includegraphics[width=0.96\linewidth]{figures/open_problem_crop.png}
\end{center}

Thus it suffices to find one admissible \(X\), with \(c=1\), for which the
normalized moment
\[
 M(p):=\frac{Y_p(0)}{\Gamma(p+1)}
\]
is not nonincreasing on \((-1,0)\).

\section{Proof intuition}

For the exponential extremizer \(X(t)=e^{-t}\), every normalized moment is
exactly one.  We perturb it by making \(e^tX(t)\) decrease smoothly only after
time one.  This preserves both differential hypotheses and avoids the
singularity at the origin.  Integration by parts then shows that every
negative normalized moment is \emph{strictly larger} than one.  Nevertheless,
the Gamma factor diverges as \(p\downarrow-1\), forcing these normalized
moments back to one at the left endpoint.  A nonincreasing function cannot
have that behavior: starting at the limiting value one, it cannot later rise
strictly above one.

\section{Construction and counterexample}

Put \(c=1\), and define a continuously differentiable, nondecreasing function
\(G:[0,\infty)\to[0,1/2]\) by
\[
G(t)=
\begin{cases}
0, & 0\leq t\leq1,\\[2pt]
\dfrac12\bigl(3(t-1)^2-2(t-1)^3\bigr), & 1<t<2,\\[4pt]
\dfrac12, & t\geq2.
\end{cases}                                                   \tag{2}
\]
Set
\[
 X(t)=e^{-t}\bigl(1-G(t)\bigr).                              \tag{3}
\]

\begin{theorem}
The function \(X\) in \emph{(3)} is positive and \(C^1\), satisfies both
hypotheses of Theorem~2.1, and has finite \(Y_p(0)\) for every \(p>-1\).
There exist \(-1<p_1<p_2<0\) for which
\[
 \frac{Y_{p_2}(0)}{\Gamma(p_2+1)}
 >
 \frac{Y_{p_1}(0)}{\Gamma(p_1+1)}.
\]
Consequently, the extension proposed in Remark~2.2 is false.
\end{theorem}

\begin{proof}
The cubic in (2) has derivative \(3(t-1)(2-t)\) before multiplication
by \(1/2\).  Hence \(G\in C^1\), \(G'\geq0\), and \(0\leq G\leq1/2\).
It follows directly that
\[
 X'(t)=-e^{-t}\bigl(1-G(t)+G'(t)\bigr)<0,
 \qquad
 \frac{d}{dt}\bigl(e^tX(t)\bigr)=-G'(t)\leq0.               \tag{4}
\]
In particular, all the source hypotheses hold.  Near zero, \(-X'(t)=e^{-t}\),
while at infinity \(-X'(t)=\tfrac12e^{-t}\), so
\(Y_p(0)=\int_0^\infty t^p(-X'(t))\,dt\) is finite for every \(p>-1\).

For such \(p\), ordinary integration by parts is legitimate because \(G\)
vanishes on \([0,1]\) and is bounded.  From (4),
\begin{align*}
Y_p(0)
 &=\int_0^\infty t^pe^{-t}\bigl(1-G(t)+G'(t)\bigr)\,dt\\
 &=\Gamma(p+1)-p\int_0^\infty G(t)t^{p-1}e^{-t}\,dt.         \tag{5}
\end{align*}
Indeed, the boundary term in integrating the \(G'\)-term is zero, and its
\(G(t)t^pe^{-t}\) contribution cancels the preceding \(-G\)-term.

Write
\[
 J(p)=\int_0^\infty G(t)t^{p-1}e^{-t}\,dt.
\]
Because \(G=0\) on \([0,1]\), the function \(J\) is finite and continuous
even on a neighborhood of \([-1,0]\).  Moreover \(J(p)>0\).  Thus (5) gives
\[
 M(p)=1+\frac{(-p)J(p)}{\Gamma(p+1)}>1
 \qquad(-1<p<0).                                             \tag{6}
\]
On the other hand, \(J(p)\to J(-1)<\infty\), while
\(\Gamma(p+1)\to\infty\), as \(p\downarrow-1\).  Therefore
\[
 \lim_{p\downarrow-1}M(p)=1.                                \tag{7}
\]
Fix any \(p_2\in(-1,0)\), for instance \(p_2=-1/2\).  By (6),
\(M(p_2)>1\); by (7), there is a \(p_1\in(-1,p_2)\) sufficiently close to
\(-1\) such that \(M(p_1)<M(p_2)\).  This is the strict reverse of (1).
\end{proof}

\section{Robustness and verification}

The cubic profile is only a convenient explicit choice.  The same proof works
for every nonzero \(C^1\) nondecreasing \(G\) that vanishes on a neighborhood
of zero, is eventually constant, and satisfies \(0\leq G<1\).  Thus the
failure is a stable family of counterexamples rather than an isolated
numerical accident.

The script \texttt{code/check\_counterexample.py}, run with 80 decimal digits,
evaluates the exact right-hand side of (5).  It reports
\[
 M(-0.99)=1.0002754937621309\ldots,
 \qquad
 M(-0.5)=1.0106701090171454\ldots,
\]
and confirms the displayed strict reversal.  These decimal values are only
QA; the existence argument from (6)--(7) is the proof.

\section{Novelty check, limitations, and review focus}

A bounded search on 17 August 2026 checked the run's registry, solution,
attempt, and proof-gap indexes; arXiv queries using the paper id, author,
``negative \(p\)'', and normalized Green-function moments; OpenAlex records
and title/keyword searches; and Crossref title/author searches.  No later
paper or prior packet answering this exact negative-exponent conjecture was
found.  OpenAlex listed no citing works for the source record.  This supports,
but does not prove, novelty; the search was not an exhaustive MathSciNet or
zbMATH review.

The result addresses the one-dimensional conjecture of Remark~2.2 exactly.
It does not by itself realize (3) as a Green-function distribution, so it does
not disprove the later geometric conjectures in Remarks~3.1 and~4.1.  Human
review should focus on the sign and boundary terms in (5), and on whether
``Theorem~2.1(i) is still valid'' is intended to retain precisely the stated
one-dimensional hypotheses when \(p_1<0\).

\begin{thebibliography}{1}
\bibitem{Xiao}
Jie Xiao,
\emph{Optimal Monotonicity of \(L^p\) Integral of Conformal Invariant Green
Function}, arXiv:0708.2644 (2007), especially Theorem~2.1 and Remark~2.2.
\end{thebibliography}

\end{document}
